\documentclass[11pt,a4paper]{article}

% ── Paquetes ────────────────────────────────────────────────────────────────
\usepackage[utf8]{inputenc}
\usepackage[T1]{fontenc}
\usepackage[spanish]{babel}
\usepackage[margin=2.5cm]{geometry}
\usepackage{amsmath,amssymb}
\usepackage{graphicx}
\usepackage{booktabs}
\usepackage{xcolor}
\usepackage{hyperref}
\usepackage{todonotes}
\usepackage{enumitem}

\hypersetup{
    colorlinks=true,
    linkcolor=blue!60!black,
    citecolor=blue!60!black,
    urlcolor=blue!60!black
}

% Caja de nota de enfoque
\newcommand{\notaenfoque}[1]{%
    \begin{quote}
    \small\itshape\color{gray!70!black} #1
    \end{quote}
}

\title{Un algoritmo de emulación estadística para reducir el costo de\\
       simulaciones basadas en agentes: caso de estudio de olas de\\
       calor en adultos mayores}
\author{Matías Escudero}
\date{Borrador --- \today}

\begin{document}

\maketitle

\notaenfoque{
\textbf{Nota de enfoque} (según indicación de la profesora, reunión 31 jul 2026):
El foco del paper NO es ``un simulador de olas de calor''. Es un
\textbf{algoritmo para reducir el costo computacional de simulaciones
basadas en agentes}, aplicado a un caso de estudio de exposición al calor
en adultos mayores. La estructura sigue exactamente esa lógica: (1) modelar
y validar el caso de estudio, (2) mostrar cómo el algoritmo propuesto
(emulador PCA+GP) rinde contra alternativas en tiempo de ejecución y calidad.

\textbf{Actualización (reunión 14 ago 2026):} escribir directamente en
\textbf{inglés} de aquí en adelante (no traducir al final) --- la versión en
español de la tesis se deriva después, más extendida. Este borrador se deja
en español por ahora solo porque así estaba empezado; migrar antes de la
próxima reunión (28 ago).
}

\begin{abstract}
Las simulaciones basadas en agentes (ABM) espaciales son una herramienta
clave para estudiar fenómenos de salud pública urbana, pero su costo
computacional limita su uso en aplicaciones que requieren respuesta casi
inmediata --- por ejemplo, sistemas de alerta temprana ante olas de calor.
Presentamos un algoritmo de emulación estadística (diseño Latin Hypercube +
compresión por componentes principales + Gaussian Process) que reduce el
tiempo de una consulta de $\sim$65--100 segundos a $\sim$2,3 milisegundos
(una aceleración de $\sim$28.000--44.000$\times$), manteniendo una fidelidad
de 87,5\% (\emph{skill}) frente al modelo original, validada mediante
\emph{leave-one-out}. Aplicamos el algoritmo a un caso de estudio de
exposición al calor en adultos mayores (65+) en el Gran Santiago, Chile
(34 comunas, 1,19 millones de adultos mayores), comparando el proceso
Gaussiano contra Random Forest, interpolación RBF y \emph{Polynomial Chaos
Expansion}. El proceso Gaussiano obtiene la mejor precisión de las cuatro
alternativas y es el único que entrega incertidumbre puntual calibrada
(92\% de cobertura del intervalo de confianza al 95\%), una propiedad
indispensable para sistemas de decisión en tiempo real.
\end{abstract}

% ══════════════════════════════════════════════════════════════════════════
\section{Introducción}
% ══════════════════════════════════════════════════════════════════════════

\begin{itemize}[leftmargin=*]
    \item Motivación: olas de calor y mortalidad en adultos mayores ---
        contexto real reciente (ej. evento de calor en París 2026,
        $\sim$85\% de las muertes en adultos mayores --- usar como gancho,
        con cita de prensa/fuente oficial).
    \item El problema computacional: los ABM espaciales de exposición son
        caros de correr (en nuestro caso, 65--100 s por escenario de 3 días
        en Gran Santiago), lo que impide su uso en un sistema operacional
        que deba responder a un pronóstico meteorológico en tiempo real.
    \item La propuesta: un algoritmo de emulación (\emph{surrogate model})
        que aprende la relación entre condiciones meteorológicas y el mapa
        de riesgo resultante, sin volver a correr el ABM por cada consulta.
    \item Contribución dual:
    \begin{enumerate}
        \item Un algoritmo general de reducción de costo para ABM espaciales
            de salida de alta dimensión (aplicable más allá de este caso de
            estudio).
        \item Su aplicación y validación en un caso de estudio real y
            actual: exposición al calor en adultos mayores, Gran Santiago.
    \end{enumerate}
\end{itemize}

% ══════════════════════════════════════════════════════════════════════════
\section{Trabajo relacionado}
% ══════════════════════════════════════════════════════════════════════════

\paragraph{Emulación de simuladores costosos (fundamento metodológico).}
Sacks et al. (1989), \emph{Design and Analysis of Computer Experiments}
\todo[inline,color=orange!40]{Verificar cita exacta}; Kennedy \& O'Hagan
(2001), \emph{Bayesian calibration of computer models}
\todo[inline,color=orange!40]{Verificar cita exacta}; Higdon et al. (2008),
\emph{Computer Model Calibration Using High-Dimensional Output}, JASA ---
base directa del pipeline PCA+GP para salida espacial de alta dimensión;
Gramacy (2020), \emph{Surrogates: Gaussian Process Modeling, Design, and
Optimization}.

\paragraph{Emulación de ABM específicamente (precedente directo).}
\emph{Towards real-time predictions using emulators of agent-based models},
Journal of Simulation (2024); \emph{Combining an agent-based model with
Gaussian process emulation} (variante Ómicron, Noruega), Epidemics (2026);
\emph{Gaussian process emulation of spatio-temporal outputs of a 2D inland
flood model}, Water Research (2022).

\paragraph{Comparación GP vs. otras técnicas de emulación (justifica la elección).}
\emph{A comparison of Gaussian processes and polynomial chaos emulators},
Phil. Trans. Royal Society A (2025); Owen et al., \emph{A comparison of
polynomial chaos and Gaussian process emulation for UQ in computer
experiments}; \emph{Emulation of CPU-demanding reactive transport models: GP
vs PCE vs redes neuronales}.

\paragraph{Extensiones / trabajo futuro citable.}
\emph{Gaussian Process Latent Factor Regression for Low-Data Problems with
High-Dimensional Outputs} (2026); \emph{Multivariate Gaussian Process
Emulators With Nonseparable Covariance Structures} (coregionalización).

\paragraph{Dominio de aplicación.}
Vulnerabilidad al calor urbano --- el vacío que este trabajo llena,
combinando ABM + emulador + tiempo real: \emph{Urban-Hazard Risk Analysis:
Heat-Related Risks in the Elderly in Major Italian Cities}; \emph{A Novel
Urban Heat Vulnerability Analysis: ML + Remote Sensing}, Remote Sensing
(2024).

% ══════════════════════════════════════════════════════════════════════════
\section{Caso de estudio: modelo de exposición al calor}
% ══════════════════════════════════════════════════════════════════════════

\subsection{Modelo de agentes}

Población: 1 agente representa a 10 adultos mayores de la misma manzana
censal (comparten vulnerabilidad/ubicación, difieren en decisión de
salir/propósito/horario, muestreado desde la Encuesta Origen-Destino Gran
Santiago 2012, filtrada a 60+).

Arquitectura de dos capas:
\begin{itemize}[leftmargin=*]
    \item \textbf{Capa invariante} (no depende del clima): red peatonal,
        rutas peatonales deduplicadas por par origen-destino, NDVI a lo
        largo de cada ruta. Se calcula una vez por zona urbana.
    \item \textbf{Capa de escenario} (depende del clima): modelo WBGT
        aditivo aplicado sobre la capa invariante.
\end{itemize}

Modelo de calor --- WBGT aditivo (no multiplicativo):
\begin{align}
    \mathrm{WBGT}_{\mathrm{eff}} &= \mathrm{WBGT}_{\mathrm{ambiente}} -
        \mathrm{ndvi\_cooling} \\
    \mathrm{ndvi\_cooling} &= \alpha_{\mathrm{ndvi}} \cdot
        \mathrm{ndvi\_norm}(\mathrm{ruta}), \quad
        \alpha_{\mathrm{ndvi}} = 2.5\,^\circ\mathrm{C}
        \text{ (Bowler et al. 2010)} \\
    \mathrm{wbgt\_umbral}_{\mathrm{agente}} &= 27\,^\circ\mathrm{C}
        \text{ (NIOSH 2016)} - \delta(\mathrm{vuln\_group})
\end{align}
con $\delta = \{\text{baja}: 0,\ \text{media}: 1.5,\ \text{alta}: 3.0\}$
(Kenney \& Munce 2003). Acumulación multi-día con
\texttt{nocturnal\_decay} $= 0.80$.

\subsection{Escala del caso de estudio}

\begin{table}[h]
\centering
\begin{tabular}{lrr}
\toprule
 & Peñalolén (piloto) & Gran Santiago (caso final) \\
\midrule
Comunas & 1 & 34 \\
Manzanas censales & $\sim$1.600 & 42.800 \\
Nodos de red peatonal & 12.535 & 212.314 \\
Adultos mayores 65+ & --- & 1,19 millones \\
Tiempo capa invariante (1 vez) & --- & 13,5 min \\
Tiempo por escenario (3 días) & --- & 65--100 s \\
\bottomrule
\end{tabular}
\caption{Escala del caso de estudio, Peñalolén vs. Gran Santiago.}
\label{tab:escala}
\end{table}

% ══════════════════════════════════════════════════════════════════════════
\section{Validación del caso de estudio}
% ══════════════════════════════════════════════════════════════════════════

\subsection{Validación temporal (RM completa) --- resultado fuerte}

Usando datos reales de defunciones DEIS (RM, 65+, 2005--2023): el exceso de
mortalidad durante olas de calor documentadas correlaciona con la
\textbf{duración} de la ola más que con su intensidad pico. Regresión
múltiple estandarizada (exceso acumulado $\sim$ Tmax\_max + n\_dias):
$R^2=0{,}55$; la duración explica $\sim$96\% de la varianza relativa; el
Tmax pico no es significativo ($p=0{,}695$).

El mismo patrón se replica de forma independiente en egresos hospitalarios
(N17 insuficiencia renal aguda + E87 trastorno electrolítico, RM 60+): la
tasa correlaciona fuertemente con la Tmax media/p90 del verano
($r=0{,}84$--$0{,}88$, $p<0{,}05$), no con el pico ($r=0{,}52$, n.s.).

\textit{Este hallazgo fue corroborado independientemente por un
epidemiólogo experto (ex-Director de Epidemiología, MINSAL) como
consistente con literatura documentada sobre duración de eventos de calor.}

\subsection{Validación espacial (comuna / establecimiento) --- resultado nulo, documentado con rigor}

Se probó si el riesgo simulado (ABM y también el índice de vulnerabilidad
estático) correlaciona espacialmente con datos reales de salud:

\begin{itemize}[leftmargin=*]
    \item \textbf{Nivel comuna} ($n=34$): mortalidad por causas sensibles al
        calor (5 eventos reales con meteorología día a día vía
        ERA5/Open-Meteo) y egresos N17+E87 anuales --- sin correlación
        significativa en ningún caso, incluso corrigiendo un artefacto de
        \emph{pooling} entre eventos (naive $\rho=0{,}31$, $p<0{,}001$
        desaparece a $\rho=0{,}02$, n.s. al estandarizar dentro de evento).
    \item \textbf{Nivel establecimiento de salud} ($n=25$): crosswalk
        georreferenciado entre códigos de urgencias DEIS y catastro nacional
        MINSAL (93--100\% de cobertura validada), \emph{catchment} por
        cercanía, urgencias por causa cardiovascular específica, con
        \emph{baseline} estacionalmente correcto --- tampoco significativo.
\end{itemize}

\textbf{Diagnóstico}: no es una falla del modelo --- el gradiente de
vulnerabilidad simulado es correcto y grande (comunas de alto ingreso
$\sim$0,2 vs.\ vulnerables $\sim$10 en \texttt{heat\_mean}, $\sim$65$\times$
de rango, geográficamente donde corresponde). Es el \textbf{problema de área
pequeña / falacia ecológica}: con conteos de salud pequeños por unidad
espacial, el ruido de Poisson domina la señal real --- confirmado
independientemente por el experto consultado.

\textbf{Trabajo en curso}: solicitud vía Ley de Transparencia de egresos
agregados por establecimiento$\times$día$\times$grupo\_edad$\times$diagnóstico
(siguiendo recomendación directa del experto), para repetir la validación a
resolución más fina.

\subsection{Validación por experto (planificada)}

Siguiendo indicación de la profesora: formalizar un proceso de validación
contra experto mediante un instrumento estructurado --- generar corridas del
simulador para escenarios conocidos, mostrar los resultados, y recoger su
evaluación formal basada en experiencia profesional en epidemiología del
calor.

\subsection{Métricas de evaluación del simulador}

Métricas actualmente implementadas:
\begin{enumerate}[leftmargin=*]
    \item Número de personas afectadas (agregado, KPI poblacional).
    \item Riesgo por persona (\texttt{heat\_load} individual, variable en el tiempo).
    \item Riesgo promedio por zona (agregación espacial, variable en el tiempo).
\end{enumerate}

\todo[inline,color=orange!40]{Pendiente: revisión de literatura sobre
métricas adicionales usadas en estudios/simulaciones de impacto de olas de
calor (pedido explícito de la profesora, 31 jul 2026).}

% ══════════════════════════════════════════════════════════════════════════
\section{El algoritmo de emulación}
% ══════════════════════════════════════════════════════════════════════════

\subsection{Diseño}

\emph{Latin Hypercube Sampling} (40 puntos, $T_{\max} \in [26,38]\,^\circ$C,
$HR_{\min} \in [15,50]$\%, $Rs_{\max} \in [700,1050]$ W/m$^2$) $\to$ correr
el ABM en cada punto $\to$ PCA sobre los 40 mapas resultantes ($K=12$
componentes, 98,9\% de varianza explicada) $\to$ un Gaussian Process por
componente (\emph{kernel} \texttt{ConstantKernel} $\times$ \texttt{RBF(ARD)}
$+$ \texttt{WhiteKernel}) $\to$ predicción vía GP + PCA inversa.

\subsection{Benchmark: GP vs.\ alternativas}

Validación cruzada \emph{leave-one-out} completa (PCA re-ajustada en cada
\emph{fold}, sin fuga de información), sobre el banco real de 40 escenarios,
Gran Santiago:

\begin{table}[h]
\centering
\begin{tabular}{lrrrrrl}
\toprule
Método & RMSE & Skill & Cobertura IC95\% & Entrenar & Predecir & Incertidumbre \\
\midrule
\textbf{GP} & \textbf{1,50} & \textbf{87,5\%} & \textbf{92\%} & 810 ms & 2,5 ms & nativa, local \\
Random Forest & 2,75 & 77,1\% & 99\%$^*$ & 2.328 ms & 131 ms & proxy, no calibrada \\
RBF & 1,79 & 85,1\% & N/A & 0,3 ms & 0,04 ms & ninguna \\
PCE (orden 2) & 1,82 & 84,8\% & N/A & 88 ms & 36 ms & global, no puntual \\
\bottomrule
\end{tabular}
\caption{Benchmark GP vs.\ alternativas, \emph{leave-one-out} $n=40$, Gran
Santiago. $^*$Cobertura de Random Forest artificialmente alta --- dispersión
entre árboles, no incertidumbre probabilística calibrada.}
\label{tab:benchmark}
\end{table}

GP obtiene la mejor precisión de las cuatro alternativas \textbf{y} es el
único con incertidumbre puntual bien calibrada --- propiedad indispensable
para un sistema operacional que debe comunicar cuánto confiar en cada
predicción.

\todo[inline,color=orange!40]{Benchmark preliminar con $n=40$; repetir sobre
el banco completo de $\sim$300 escenarios antes de la versión final.}

\subsection{Tiempo de ejecución aislado}

\begin{table}[h]
\centering
\begin{tabular}{lr}
\toprule
Componente & Tiempo \\
\midrule
GP predict (12 componentes, media + desviación) & 2,02 ms \\
PCA inversa (12 pesos $\to$ mapa de 42.800 manzanas) & 0,27 ms \\
\textbf{Total} & \textbf{2,29 ms} \\
\midrule
\emph{(referencia)} 1 corrida del ABM (Gran Santiago, 3 días) & 65.000--100.000 ms \\
\bottomrule
\end{tabular}
\caption{Tiempo de ejecución aislado del emulador vs.\ el ABM completo.
Aceleración: $\sim$28.400$\times$--43.700$\times$.}
\label{tab:tiempos}
\end{table}

\subsection{Escalamiento: tiempo de ejecución vs.\ número de agentes}

Pedido explícito de la profesora (reunión 14 ago 2026): un gráfico con eje
$X$ = número de agentes simulados, eje $Y$ = tiempo de ejecución, comparando
el ABM completo contra el emulador GP (línea plana, $\sim$2,3 ms
independiente de $N$, porque el costo del ABM ya quedó ``destilado'' en el
banco de entrenamiento durante el diseño LHS).

Resultado medido (Gran Santiago, subconjuntos de 10--100\% de los 45.047
\emph{walkers}, un día de simulación por punto):

\begin{table}[h]
\centering
\begin{tabular}{rrr}
\toprule
N.\ de agentes & Tiempo ABM (1 día) & Tiempo emulador GP \\
\midrule
4.504 & 3,60 s & 2,29 ms \\
11.261 & 5,08 s & 2,29 ms \\
22.523 & 8,03 s & 2,29 ms \\
33.785 & 11,11 s & 2,29 ms \\
45.047 & 13,73 s & 2,29 ms \\
\bottomrule
\end{tabular}
\caption{Tiempo de ejecución del ABM vs.\ número de agentes (1 día de
simulación), comparado contra el tiempo constante del emulador GP.}
\label{tab:tiempo_vs_agentes}
\end{table}

El tiempo del ABM crece de forma aproximadamente lineal con el número de
agentes ($t \approx 2{,}5\text{s} + 0{,}25\text{ms} \times N_{\mathrm{agentes}}$,
ajuste por mínimos cuadrados sobre los 5 puntos), reflejando un costo fijo de
inicialización más un costo marginal por agente. El emulador GP, en cambio,
es \textbf{constante} respecto a $N$ --- su costo escala con la
\emph{compresión} del mapa (12 componentes PCA), no con el número de agentes
que generaron ese mapa. Esta es la propiedad clave: mientras más grande el
caso de estudio (más agentes, más ciudad), mayor la ventaja relativa del
emulador.

\todo[inline,color=green!30!black]{Falta: figura (gráfico de líneas, eje $Y$
en escala logarítmica dado el contraste de 4 órdenes de magnitud) --- datos
ya en \texttt{scenarios/gran\_santiago/tiempo\_vs\_agentes.csv}. Nota: esta
corrida usa 1 día de simulación por agilidad; la cifra de referencia de
65--100s (Tabla~\ref{tab:tiempos}) es para 3 días --- homologar antes de la
versión final.}

\subsection{Estudio de ablación}

Pedido explícito de la profesora (reunión 14 ago 2026): un estudio de
\emph{ablación} --- remover/agregar sistemáticamente componentes del modelo
y medir el impacto de cada uno en el resultado final, para verificar que
ningún componente es redundante y cuantificar el aporte real de cada parte.

\todo[inline,color=orange!40]{PENDIENTE --- diseñar y correr. Candidatos a
ablacionar (propuesta inicial, a confirmar):
\begin{itemize}[leftmargin=*,nosep]
    \item Enfriamiento por NDVI (\texttt{ndvi\_alpha}: activo vs.\ 0)
    \item Diferenciación de umbral WBGT por grupo de vulnerabilidad
        (\texttt{vuln\_delta}: activo vs.\ uniforme)
    \item Velocidad de caminata diferenciada por vulnerabilidad (activo vs.
        velocidad única)
    \item Decaimiento nocturno multi-día (\texttt{nocturnal\_decay}: 0,80
        vs.\ 0, es decir sin memoria entre días)
\end{itemize}
Para cada configuración: correr el mismo banco de escenarios, medir el
cambio en el mapa de riesgo resultante (skill/RMSE respecto a la
configuración completa) --- esto muestra qué tan sensible es el resultado a
cada mecanismo del modelo.}

\subsection{Pruebas estadísticas adicionales}

\todo[inline,color=orange!40]{PENDIENTE (pedido de la profesora): agregar
prueba estadística formal (ej.\ test $z$) sobre la diferencia de skill entre
GP y las alternativas del benchmark (Tabla~\ref{tab:benchmark}), no solo
reportar el punto estimado.}

% ══════════════════════════════════════════════════════════════════════════
\section{Resultados}
\todo[inline,color=orange!40]{Pendiente --- se completa con la corrida del
banco completo y la validación por experto.}
% ══════════════════════════════════════════════════════════════════════════

% ══════════════════════════════════════════════════════════════════════════
\section{Análisis de escenarios de mitigación}
\label{sec:mitigacion}
% ══════════════════════════════════════════════════════════════════════════

\todo[inline,color=purple!40]{SECCIÓN NUEVA (reunión 14 ago 2026) --- según
la profesora, este debería ser un \textbf{capítulo completo de la tesis},
aquí solo un esbozo para el paper. Usa directamente la capacidad de
``reciclaje de escenarios'' del emulador: el mapa de riesgo identifica
\emph{dónde} conviene intervenir, y el emulador permite estimar el efecto de
una intervención en milisegundos, sin recorrer todo el banco de nuevo.}

Estrategias de mitigación de calor urbano documentadas en otras ciudades,
como referencia para el diseño de escenarios \emph{what-if}:

\begin{itemize}[leftmargin=*]
    \item \textbf{Nebulización en altura} (provincias en China): sistemas en
        azoteas que emiten micro-gotas de agua a alta presión, enfriando el
        aire por evaporación --- reducción reportada de 5--8$\,^\circ$C.
    \item \textbf{Refrigeración radiativa} (Hong Kong): materiales que
        expulsan calor de edificios hacia el exterior mediante radiación
        infrarroja.
    \item \textbf{Estructuras de sombra retráctiles} (Corea, Shanghai):
        sobre calles peatonales, desplegadas en verano, guardadas en
        invierno --- generan pasillos de sombra continua.
    \item \textbf{Techos blancos y arquitectura de sombra}: pórticos y
        orientación estratégica de plazas para maximizar sombra natural.
    \item Puntos de hidratación en zonas de alto tránsito peatonal de
        adultos mayores.
\end{itemize}

\textbf{Diseño del análisis propuesto}: usando el mapa de riesgo simulado
(las zonas de mayor \texttt{heat\_load} acumulado, ``zonas rojas''),
proponer ubicaciones candidatas para cada intervención, aplicar su efecto
conocido (ej.\ $-5\,^\circ$C localizado $\to$ reducción de
$\mathrm{WBGT}_{\mathrm{eff}}$ en las rutas que cruzan esa zona) y volver a
consultar el emulador para estimar: (i) cuántas zonas dejan de estar en
riesgo alto, (ii) reducción del riesgo total agregado, (iii) comparación de
costo-efectividad entre estrategias. Esto conecta directamente con la
propuesta original de tesis (sección de aplicabilidad para toma de
decisiones) y es un caso de uso concreto de \emph{scenario recycling}: cada
consulta ``qué pasa si pongo sombra acá'' se resuelve en milisegundos, no
recorriendo el ABM completo.

% ══════════════════════════════════════════════════════════════════════════
\section{Discusión}
\todo[inline,color=orange!40]{Pendiente.}
% ══════════════════════════════════════════════════════════════════════════

% ══════════════════════════════════════════════════════════════════════════
\section{Conclusión y trabajo futuro}
% ══════════════════════════════════════════════════════════════════════════

\begin{itemize}[leftmargin=*]
    \item Banco completo ($\sim$300 escenarios) vía SLURM \emph{job arrays}
        en el cluster.
    \item GP multi-salida con coregionalización (en vez de 12 GPs
        independientes) --- ver \emph{Gaussian Process Latent Factor
        Regression for Low-Data Problems} (2026).
    \item Validación espacial a resolución de establecimiento con datos de
        fecha exacta (pendiente respuesta de solicitud de Transparencia).
    \item Validación formal contra experto.
    \item API + visor para consulta en tiempo real de pronósticos DMC.
\end{itemize}

% ══════════════════════════════════════════════════════════════════════════
\section*{Notas internas --- checklist para la próxima reunión (28 ago)}
{\small
\textbf{De la reunión 31 jul:}
\begin{itemize}[leftmargin=*]
    \item[$\square$] Buscar métricas adicionales de impacto de olas de calor
        en literatura (pedido directo de la profesora)
    \item[$\square$] Confirmar citas exactas de Sacks (1989) y Kennedy \&
        O'Hagan (2001)
    \item[$\square$] Avanzar validación por urgencias con los tips de
        Cristian
    \item[$\square$] Enviar solicitud de Transparencia --- \textbf{aún sin
        enviar según reunión 14 ago}
    \item[$\square$] Diseñar el instrumento de validación por experto
        (formulario + corridas de ejemplo)
\end{itemize}

\textbf{Nuevo de la reunión 14 ago:}
\begin{itemize}[leftmargin=*]
    \item[$\square$] \textbf{Migrar el paper a inglés} (escribir directo en
        inglés de ahora en adelante)
    \item[$\square$] Diseñar y correr el \textbf{estudio de ablación}
        (ver sección de ablación --- candidatos ya propuestos)
    \item[$\boxminus$] Gráfico tiempo de ejecución vs.\ número de agentes ---
        \textbf{datos ya obtenidos hoy}, falta la figura
    \item[$\square$] Prueba estadística formal (test $z$ o similar) sobre la
        diferencia de skill del benchmark GP vs.\ alternativas
    \item[$\square$] Esbozar el capítulo de \textbf{análisis de escenarios de
        mitigación} (sección nueva agregada --- shade structures,
        nebulización, puntos de hidratación; usar el emulador para estimar
        efecto de cada intervención)
    \item[$\square$] Evaluar si medir consumo de RAM como métrica adicional
        (sugerencia de la profesora --- probablemente no relevante a esta
        escala, pero revisar si crece con más agentes)
    \item[$\square$] Revisar si el estudio de Cristian García (nivel
        regional, olas de calor) es citable --- confirmó conclusiones
        macro alineadas con las nuestras
    \item[$\square$] Completar secciones de Resultados y Discusión
\end{itemize}
}

\end{document}
